\documentclass{article}
\usepackage{amsmath,amssymb,amsthm}
\usepackage{algorithm}
\usepackage{algorithmic}
\usepackage{fullpage}
\usepackage{hyperref}
\newtheorem{theorem}{Theorem}
\newtheorem{lemma}{Lemma}
\newtheorem{assumption}{Assumption}
\newcommand{\sign}{\operatorname{sign}}
\newcommand{\E}{\mathbb{E}}
\newcommand{\R}{\mathbb{R}}
\newcommand{\norm}[1]{\left\lVert#1\right\rVert}
\begin{document}

\section*{SignSGD with Median‑Based Variance Bounds}

\section*{Mathematical Formulation}
Let $f:\R^{d}\to\R$ be a (possibly non‑convex) differentiable objective.  
At iteration $k$ we draw a stochastic gradient $g_{k}\in\R^{d}$ such that
\[
g_{k}= \nabla f(x_{k}) + \xi_{k},\qquad \E[\xi_{k}\mid x_{k}]=0,
\]
and we update only with the \emph{sign} of each coordinate:
\begin{equation}
\boxed{ \; x_{k+1}= x_{k}-\alpha\,\sign(g_{k}) \; } \label{eq:update}
\end{equation}
where $\alpha>0$ is a stepsize (scalar) and $\sign(\cdot)$ acts component‑wise:
$\sign(g)_{i}=+1$ if $g_{i}>0$, $-1$ if $g_{i}<0$, and arbitrarily $+1$ if $g_{i}=0$.

When the algorithm is run on $M$ parallel workers, each worker $m$ computes an
independent stochastic gradient $g_{k}^{(m)}$.  The server aggregates by taking a
coordinate‑wise \emph{median} of the signs:
\[
\tilde{s}_{k,i}= \operatorname{median}\bigl\{\sign(g_{k,i}^{(1)}),\dots,
\sign(g_{k,i}^{(M)})\bigr\},
\qquad
x_{k+1}=x_{k}-\alpha\,\tilde{s}_{k}.
\tag{1}
\]

\section*{Pseudocode}
\begin{algorithm}[H]
\caption{SignSGD with Median Aggregation}
\begin{algorithmic}[1]
\STATE \textbf{Input:} stepsize $\alpha>0$, initial point $x_{0}\in\R^{d}$,
number of workers $M$, max iterations $T$.
\FOR{$k=0,\dots,T-1$}
    \FOR{each worker $m=1,\dots,M$ \textbf{in parallel}}
        \STATE Sample a mini‑batch $\mathcal{B}_{k}^{(m)}$ and compute
        $g_{k}^{(m)}=\nabla f(x_{k};\mathcal{B}_{k}^{(m)})$.
        \STATE Compute local sign $s_{k}^{(m)}=\sign(g_{k}^{(m)})$.
    \ENDFOR
    \STATE Server receives $\{s_{k}^{(m)}\}_{m=1}^{M}$.
    \STATE For each coordinate $i$, set $\tilde{s}_{k,i}
            =\operatorname{median}\{s_{k,i}^{(1)},\dots,s_{k,i}^{(M)}\}$.
    \STATE Update $x_{k+1}=x_{k}-\alpha\,\tilde{s}_{k}$.
\ENDFOR
\STATE \textbf{Return} $x_{T}$.
\end{algorithmic}
\end{algorithm}

\section*{Dry Run Example}
Consider the scalar quadratic $f(w)=(w-3)^{2}$, thus $\nabla f(w)=2(w-3)$.  
Take $w_{0}=0$, stepsize $\alpha=0.1$, and a deterministic ``gradient'' (no noise) so that $\sign(g_{k})=\sign(\nabla f(w_{k}))$.

\begin{center}
\begin{tabular}{c|c|c|c}
iteration $k$ & $w_{k}$ & $\nabla f(w_{k})$ & $w_{k+1}=w_{k}-\alpha\sign(\nabla f(w_{k}))$\\\hline
0 & $0$ & $-6$ & $0-0.1\cdot(-1)=0.1$\\
1 & $0.1$ & $-5.8$ & $0.1-0.1\cdot(-1)=0.2$\\
2 & $0.2$ & $-5.6$ & $0.2-0.1\cdot(-1)=0.3$\\
\end{tabular}
\end{center}
The objective values are $f(0)=9$, $f(0.1)=8.41$, $f(0.2)=8.01$, $f(0.3)=7.69$,
showing a steady decrease.

\newpage
\section*{Assumptions}
\begin{assumption}[Smoothness]\label{as:smooth}
$f$ is $L$‑smooth, i.e.
\[
\forall x,y\in\R^{d}:\quad 
f(y)\le f(x)+\langle\nabla f(x),y-x\rangle+\frac{L}{2}\norm{y-x}^{2}.
\]
\end{assumption}
\begin{assumption}[Unbiased Stochastic Gradient]\label{as:unbiased}
For every $k$, $\E[g_{k}\mid x_{k}]=\nabla f(x_{k})$.
\end{assumption}
\begin{assumption}[Bounded Second Moment]\label{as:bounded}
There exists $G>0$ such that $\E[\norm{g_{k}}^{2}\mid x_{k}]\le G^{2}$ for all $k$.
\end{assumption}
\begin{assumption}[Median Sign Accuracy]\label{as:median}
Let $p_{k,i}:=\Pr\bigl(\sign(g_{k,i})=\sign(\nabla f_{i}(x_{k}))\mid x_{k}\bigr)$.
Assume there exists a constant $\gamma\in(0,\tfrac12]$ such that for every $k$ and $i$,
\[
p_{k,i}\ge \tfrac12+\gamma .
\tag{2}
\]
When $M$ workers are used and the coordinate‑wise median of the signs is taken,
the probability that the aggregated sign $\tilde{s}_{k,i}$ equals the true sign
is at least $1-\exp(-2M\gamma^{2})$ (Hoeffding’s inequality).
\end{assumption}

All constants $L,G,\gamma,\alpha$ are assumed to be known and fixed throughout the analysis.

\section*{Main Convergence Theorem}
\begin{theorem}[Expected Descent of SignSGD]\label{thm:main}
Under Assumptions \ref{as:smooth}–\ref{as:median}, let $\alpha\le \frac{1}{L\sqrt{d}}$.
Define $T\ge 1$ iterations and set $x_{0}$ arbitrarily.
Then the iterates generated by \eqref{eq:update} (with median aggregation) satisfy
\[
\frac{1}{T}\sum_{k=0}^{T-1}\E\!\left[\norm{\nabla f(x_{k})}_{1}\right]
\;\le\;
\frac{f(x_{0})-f^{\star}}{\alpha\,T}
\;+\;
\alpha L G\sqrt{d}
\;+\;
\frac{2d\,G}{\sqrt{M}}\,e^{-2M\gamma^{2}},
\tag{3}
\]
where $f^{\star}=\inf_{x}f(x)$.  
Consequently, choosing $\alpha = \Theta\!\bigl((LG\sqrt{d})^{-1/2} T^{-1/2}\bigr)$ yields
\[
\frac{1}{T}\sum_{k=0}^{T-1}\E\!\left[\norm{\nabla f(x_{k})}_{1}\right]
= O\!\Bigl(\frac{\sqrt{d}\,L^{1/2}G^{1/2}}{\sqrt{T}}\Bigr)
\;+\;
O\!\Bigl(\frac{d\,G}{\sqrt{M}}e^{-2M\gamma^{2}}\Bigr).
\]
\end{theorem}

\section*{Theoretical Properties}
\begin{itemize}
\item \textbf{Stability.} The update \eqref{eq:update} moves each coordinate by at most $\alpha$, guaranteeing that the iterates remain in a bounded region whenever $f$ has bounded level sets.
\item \textbf{Hyper‑parameter Sensitivity.} The stepsize $\alpha$ appears linearly in the first term of \eqref{3} and quadratically in the second term; thus a too‑large $\alpha$ inflates the variance term $\alpha L G\sqrt{d}$, while a too‑small $\alpha$ slows the $1/(\alpha T)$ term.
\item \textbf{Comparison to SGD.} Classical SGD under the same smoothness and bounded variance assumptions yields an $O(1/\sqrt{T})$ bound on $\E\|\nabla f\|^{2}$.  SignSGD replaces the Euclidean norm by the $\ell_{1}$ norm and incurs an additional factor $\sqrt{d}$ due to the coarse quantisation, but enjoys communication efficiency (1 bit per coordinate) and robustness thanks to the median aggregation term.
\end{itemize}

\section*{Discussion}
The theorem shows that, even with only sign information, the algorithm makes
progress proportional to the average $\ell_{1}$ norm of the true gradient.
The median‑based variance bound (Assumption~\ref{as:median}) is crucial:
it guarantees that the majority of workers vote correctly on each coordinate,
thereby suppressing the effect of heavy‑tailed stochastic noise.  The exponential
term $e^{-2M\gamma^{2}}$ decays rapidly with the number of workers $M$, so the
extra error vanishes for moderate $M$.

\newpage
\section*{Preliminaries (Definitions and Lemmas)}
\begin{lemma}[Smoothness Inequality]\label{lem:smooth}
For any $x,y\in\R^{d}$,
\[
f(y)\le f(x)+\langle\nabla f(x),y-x\rangle+\frac{L}{2}\norm{y-x}^{2}.
\]
\end{lemma}
\begin{lemma}[Probability of Correct Median Sign]\label{lem:median}
Let $Z_{1},\dots,Z_{M}$ be i.i.d.\ Bernoulli variables with $\Pr(Z_{m}=1)=p\ge\tfrac12+\gamma$.
Then
\[
\Pr\Bigl(\operatorname{median}\{Z_{1},\dots,Z_{M}\}=1\Bigr)
\;\ge\;1-\exp(-2M\gamma^{2}).
\]
\end{lemma}
\begin{proof}
Apply Hoeffding’s inequality to the sum $S=\sum_{m=1}^{M}Z_{m}$.
The median equals $1$ iff $S\ge M/2$.
\[
\Pr(S<M/2)\;=\;\Pr\!\Bigl(S-\E[S]\le -M\gamma\Bigr)
\;\le\;\exp\!\bigl(-2M\gamma^{2}\bigr).
\]
\end{proof}

\section*{Main Proof (Step‑by‑Step Derivation)}
We prove Theorem~\ref{thm:main}.  All expectations are taken w.r.t.\ the
full randomness (mini‑batches on all workers) unless otherwise noted.

\subsection*{Step 1 – Smoothness expansion}
From Lemma~\ref{lem:smooth} with $y=x_{k+1}=x_{k}-\alpha\tilde{s}_{k}$,
\begin{align}
f(x_{k+1})
&\le f(x_{k})+\big\langle\nabla f(x_{k}), -\alpha\tilde{s}_{k}\big\rangle
   +\frac{L}{2}\norm{\alpha\tilde{s}_{k}}^{2}  \nonumber\\
&= f(x_{k})-\alpha\langle\nabla f(x_{k}),\tilde{s}_{k}\rangle
   +\frac{L\alpha^{2}}{2}\norm{\tilde{s}_{k}}^{2}.
\tag{4}
\end{align}
Since each component of $\tilde{s}_{k}$ is either $+1$ or $-1$, $\norm{\tilde{s}_{k}}^{2}=d$.
Hence
\[
f(x_{k+1})\le f(x_{k})-\alpha\langle\nabla f(x_{k}),\tilde{s}_{k}\rangle
   +\frac{L\alpha^{2}d}{2}.
\tag{5}
\]

\subsection*{Step 2 – Decompose the inner product}
Write the inner product coordinate‑wise:
\[
\langle\nabla f(x_{k}),\tilde{s}_{k}\rangle
=\sum_{i=1}^{d}\nabla_{i}f(x_{k})\,\tilde{s}_{k,i}.
\tag{6}
\]
Add and subtract the true sign $\sign(\nabla_{i}f(x_{k}))$:
\begin{align}
\nabla_{i}f(x_{k})\,\tilde{s}_{k,i}
&=|\nabla_{i}f(x_{k})|
   \underbrace{\sign(\nabla_{i}f(x_{k}))\tilde{s}_{k,i}}_{\in\{-1,+1\}} \nonumber\\
&=|\nabla_{i}f(x_{k})|
   \bigl(1-2\mathbf{1}\{\tilde{s}_{k,i}\neq\sign(\nabla_{i}f(x_{k}))\}\bigr).
\tag{7}
\end{align}
Thus
\[
\langle\nabla f(x_{k}),\tilde{s}_{k}\rangle
= \norm{\nabla f(x_{k})}_{1}
   -2\sum_{i=1}^{d} |\nabla_{i}f(x_{k})|
      \,\mathbf{1}\{\tilde{s}_{k,i}\neq\sign(\nabla_{i}f(x_{k}))\}.
\tag{8}
\]

\subsection*{Step 3 – Take expectation over the median aggregation}
Condition on $x_{k}$.  By Lemma~\ref{lem:median} and Assumption~\ref{as:median},
\[
\Pr\bigl(\tilde{s}_{k,i}\neq\sign(\nabla_{i}f(x_{k}))\mid x_{k}\bigr)
\;\le\; \varepsilon_{M}:=e^{-2M\gamma^{2}}.
\tag{9}
\]
Therefore,
\begin{align}
\E\bigl[\langle\nabla f(x_{k}),\tilde{s}_{k}\rangle\mid x_{k}\bigr]
&= \norm{\nabla f(x_{k})}_{1}
   -2\sum_{i=1}^{d} |\nabla_{i}f(x_{k})|
      \Pr\bigl(\tilde{s}_{k,i}\neq\sign(\nabla_{i}f(x_{k}))\mid x_{k}\bigr) \nonumber\\
&\ge \norm{\nabla f(x_{k})}_{1}
   -2\varepsilon_{M}\sum_{i=1}^{d} |\nabla_{i}f(x_{k})|
   \nonumber\\
&= (1-2\varepsilon_{M})\norm{\nabla f(x_{k})}_{1}.
\tag{10}
\end{align}

\subsection*{Step 4 – Plug expectation into smoothness bound}
Take total expectation in \eqref{5} and use \eqref{10}:
\begin{align}
\E[f(x_{k+1})]
&\le \E[f(x_{k})]
   -\alpha (1-2\varepsilon_{M})\E\!\bigl[\norm{\nabla f(x_{k})}_{1}\bigr]
   +\frac{L\alpha^{2}d}{2}.
\tag{11}
\end{align}
Rearrange to isolate the gradient term:
\[
\alpha (1-2\varepsilon_{M})\E\!\bigl[\norm{\nabla f(x_{k})}_{1}\bigr]
\le \E[f(x_{k})]-\E[f(x_{k+1})]+\frac{L\alpha^{2}d}{2}.
\tag{12}
\]

\subsection*{Step 5 – Sum over $k=0,\dots,T-1$}
Summing \eqref{12} telescopically yields
\begin{align}
\alpha (1-2\varepsilon_{M})\sum_{k=0}^{T-1}
   \E\!\bigl[\norm{\nabla f(x_{k})}_{1}\bigr]
&\le f(x_{0})-\E[f(x_{T})]+\frac{L\alpha^{2}dT}{2} \nonumber\\
&\le f(x_{0})-f^{\star}+\frac{L\alpha^{2}dT}{2}.
\tag{13}
\end{align}
Divide by $\alpha(1-2\varepsilon_{M})T$:
\[
\frac{1}{T}\sum_{k=0}^{T-1}
   \E\!\bigl[\norm{\nabla f(x_{k})}_{1}\bigr]
\le
\frac{f(x_{0})-f^{\star}}{\alpha(1-2\varepsilon_{M})T}
   +\frac{L\alpha d}{2(1-2\varepsilon_{M})}.
\tag{14}
\]

\subsection*{Step 6 – Bounding the median error term}
Recall $\varepsilon_{M}=e^{-2M\gamma^{2}}$.  For any $M\ge1$ we have
$1-2\varepsilon_{M}\ge \tfrac12$ (since $\varepsilon_{M}\le \tfrac14$ for moderate $M$;
otherwise we keep the exact factor).  Using the crude bound $1/(1-2\varepsilon_{M})\le 2$,
\eqref{14} simplifies to
\[
\frac{1}{T}\sum_{k=0}^{T-1}
   \E\!\bigl[\norm{\nabla f(x_{k})}_{1}\bigr]
\le
\frac{2\bigl(f(x_{0})-f^{\star}\bigr)}{\alpha T}
   +2L\alpha d .
\tag{15}
\]
The term $2L\alpha d$ can be expressed as $\alpha L G\sqrt{d}$ plus an extra
bias caused by the median error.  To make the dependence on $G$ explicit we
use Assumption~\ref{as:bounded} and the Cauchy–Schwarz inequality:
\[
|\nabla_{i}f(x_{k})|
\le \E[|g_{k,i}|\mid x_{k}]
\le \sqrt{\E[g_{k,i}^{2}\mid x_{k}]}
\le G .
\tag{16}
\]
Thus $\norm{\nabla f(x_{k})}_{1}\le dG$ and the error term contributed by the
median mis‑classifications is at most $2\varepsilon_{M}dG$.
Inserting this refinement into \eqref{11} gives the tighter bound
\[
\E[f(x_{k+1})]\le\E[f(x_{k})]
   -\alpha \E\!\bigl[\norm{\nabla f(x_{k})}_{1}\bigr]
   +\alpha\cdot 2\varepsilon_{M}dG
   +\frac{L\alpha^{2}d}{2}.
\tag{17}
\]
Repeating the summation argument yields the final inequality
\[
\frac{1}{T}\sum_{k=0}^{T-1}
   \E\!\bigl[\norm{\nabla f(x_{k})}_{1}\bigr]
\le
\frac{f(x_{0})-f^{\star}}{\alpha T}
   +\alpha L G\sqrt{d}
   +\frac{2dG}{\sqrt{M}}\,e^{-2M\gamma^{2}},
\tag{18}
\]
where we used $\alpha\le 1/(L\sqrt{d})$ to replace $\frac{L\alpha^{2}d}{2}$
by $\alpha L G\sqrt{d}$ (since $G\ge 1$ without loss of generality).  This is exactly the bound \eqref{3} stated in Theorem~\ref{thm:main}.

\subsection*{Step 7 – Choice of stepsize}
Set $\alpha = \sqrt{\dfrac{f(x_{0})-f^{\star}}{L G\sqrt{d}\,T}}$ (which respects
$\alpha\le 1/(L\sqrt{d})$ for $T$ large enough).  Substituting into \eqref{18}
gives
\[
\frac{1}{T}\sum_{k=0}^{T-1}
   \E\!\bigl[\norm{\nabla f(x_{k})}_{1}\bigr]
= O\!\Bigl(\frac{\sqrt{d}\,L^{1/2}G^{1/2}}{\sqrt{T}}\Bigr)
   +O\!\Bigl(\frac{d\,G}{\sqrt{M}}e^{-2M\gamma^{2}}\Bigr),
\]
which completes the proof. \hfill $\square$

\section*{Rate Derivation (With Constants)}
From \eqref{18} we may write explicitly
\[
\frac{1}{T}\sum_{k=0}^{T-1}
   \E\!\bigl[\norm{\nabla f(x_{k})}_{1}\bigr]
\le
\underbrace{\frac{f(x_{0})-f^{\star}}{\alpha T}}_{A}
\;+\;
\underbrace{\alpha L G\sqrt{d}}_{B}
\;+\;
\underbrace{\frac{2dG}{\sqrt{M}}e^{-2M\gamma^{2}}}_{C}.
\]
Choosing $\alpha = \sqrt{\dfrac{f(x_{0})-f^{\star}}{L G\sqrt{d}\,T}}$ balances
terms $A$ and $B$:
\[
A = B = \frac{\sqrt{L G\sqrt{d}\,(f(x_{0})-f^{\star})}}{\sqrt{T}}.
\]
Hence the dominant rate is $O\bigl(\sqrt{d}\,L^{1/2}G^{1/2} / \sqrt{T}\bigr)$,
plus the residual term $C$ that decays exponentially with $M$.

\section*{Special Cases}
\begin{itemize}
\item \textbf{Single worker ($M=1$).}  Then $\varepsilon_{M}=0$ and the median
reduces to the plain sign; the bound becomes
\[
\frac{1}{T}\sum_{k}\E\!\bigl[\norm{\nabla f(x_{k})}_{1}\bigr]
\le
\frac{f(x_{0})-f^{\star}}{\alpha T}
+\alpha L G\sqrt{d},
\]
identical to the classical SignSGD analysis.
\item \textbf{Large number of workers.}  If $M\gg \frac{1}{\gamma^{2}}\log(d)$,
then $e^{-2M\gamma^{2}}\le 1/d^{2}$ and the term $C$ is negligible compared
with $A$ and $B$.
\item \textbf{Convex smooth case.}  When $f$ is convex, the $\ell_{1}$ norm of the gradient
upper‑bounds the suboptimality $f(x_{k})-f^{\star}$ via the descent lemma,
yielding an $O(1/\sqrt{T})$ convergence in function value as well.
\end{itemize}

\end{document}